\originalchapter{注意力与 Transformer 表征}\label{ch:attention}
\originalsection{把上下文匹配写成概率}
词表查表给同一个 token 相同的初始向量, 但“苹果很好吃”与“苹果发布产品”中的上下文不同. 逐 token 前馈网络只改变当前向量, 不能直接读取其他位置. 我们需要一个根据当前内容选择上下文, 并把取回的信息写入表征的操作. 注意力把这一任务拆为查询、键和值: 查询描述当前要匹配什么, 键描述每个候选位置怎样被匹配, 值描述选中后要带回什么. 键和值的投影可以不同, 如目录索引与记录正文的作用不同.

固定查询位置 $i$, 给候选键位置 $j$ 一个实分数 $s_{ij}$. 对允许集合 $\mathcal A_i$, 定义
\[
A_{ij}=\frac{e^{s_{ij}}}{\sum_{l\in\mathcal A_i}e^{s_{il}}}\quad(j\in\mathcal A_i),
\qquad A_{ij}=0\quad(j\notin\mathcal A_i).
\]
若允许集合非空, 这些数非负且和为 1. 可以把它们看成条件于当前查询和候选内容时, 在源位置上的一个分类概率分布. 输出 $C_i=\sum_j A_{ij}V_j$ 是在该分布下取值向量的期望. 这里的“概率”描述匹配权重, 不表示某位置有相同概率包含真答案, 也不表示模型回答正确的概率.

\begin{definition}
单头自注意力的输入 $X\in\R^{T\times d}$ 把 token 放在行上. 取 $W_Q,W_K\in\R^{d\times d_k}$, $W_V\in\R^{d\times d_v}$, 定义
\[
Q=XW_Q,\quad K=XW_K,\quad V=XW_V,\qquad
S=\frac{QK^\top}{\sqrt{d_k}}+M,\quad A=\softmax_{\rm row}(S),\quad C=AV.
\]
$Q,K$ 的形状是 $T\times d_k$, $V,C$ 是 $T\times d_v$, $S,A,M$ 是 $T\times T$. 常量掩码 $M$ 在允许位置为 0, 禁止位置记为 $-\infty$. 这个记号表达“从归一化和中排除”, 不是把无穷当成可训练参数.
\end{definition}
点积 $q_i k_j^\top$ 大, 表示在学到的查询与键坐标中相匹配; 参数训练决定哪些方向有用. 查询和值来自同一输入也不意味着它们是同一个对象. 在交叉注意力中, 查询可以来自待生成的序列, 键和值来自另一份输入或编码结果. 本章先推自注意力, 交叉注意力只需分别跟踪这两条输入来源.

缩放 $1/\sqrt{d_k}$ 有一个明确的尺度理由\cite{attention}. 假设各坐标独立, 查询与键独立, 坐标均值零、方差一, 则 $q_i k_j^\top=\sum_{r=1}^{d_k}q_{ir}k_{jr}$ 均值零、方差 $d_k$. 除以 $\sqrt{d_k}$ 后方差为 1. 点积分数过大时 softmax 可能接近只选一个键, 其大多数导数很小. 缩放改善了初始化时的尺度, 但训练后的坐标不一定满足独立同分布假设, 它也不保证每个头永不饱和.

\begin{example}
一个查询有两个允许的键, 分数为 $(0,\log3)$, 值为 $(2,0)$ 和 $(0,4)$. 求输出. 若第二键禁止, 又如何?
\end{example}
\begin{solution}
归一化分母为 $1+3=4$, 权重为 $(1/4,3/4)$, 输出是 $(1/2,3)$. 禁止第二键后它不再进入分母, 权重变为 $(1,0)$, 输出为 $(2,0)$. 不是先算 $(1/4,3/4)$ 再把第二权重改为零; 那样第一权重仍为 $1/4$, 行和不为 1, 已改变运算.
\end{solution}

稳定实现先减去每行最大允许分数, 再指数与归一化. 同减常数不改变比值, 因分子分母都乘同一个 $e^{-m_i}$. 每行至少须有一个允许位置. 全禁止行既没有概率分布, 也会在 $-\infty-(-\infty)$ 中产生未定义值. 若 padding 导致此情形, 应明确不计算该查询、给出有效的约定路径, 或拒绝输入; 随意把 NaN 改成零不等于修复了概率模型.

\originalsection{因果掩码与训练目标}
在自回归语言模型中, 输入位置 $i$ 的 logits 预测下一个 token $x_{i+1}$. 因果掩码允许 $j\le i$, 禁止 $j>i$. 允许读取当前位置 $x_i$ 不等于偷看目标, 因为目标是下一位置. 若输入与目标没有正确错开, 或注意力读到了 $x_{i+1}$, 训练损失可能很好看, 但实际生成时没有这条信息路径.

单层因果性可以直接证明. $q_i$ 只由 $X_i$ 产生, 允许的 $k_j,v_j$ 只由 $X_j,j\le i$ 产生, 所以 $C_i$ 只依赖 $X_0,\ldots,X_i$. 多层时用归纳: 若上一层位置 $j$ 的状态只依赖输入前缀至 $j$, 那么本层位置 $i$ 读取的所有 $j\le i$ 也只依赖前缀至 $i$. 逐 token 前馈、逐 token LayerNorm 和同位置残差不扩大这条信息范围. 沿全部位置做均值归一化、非因果卷积或错位数据却可能破坏它. 因而验证因果性要检查整个前向图, 不能只看到一个三角矩阵就结束.

训练时可以并行算出所有位置, 因为教师提供真实前缀, 掩码限制每个查询读取的范围. 生成时后续 token 尚未存在, 须逐步产生. 并行训练与逐步生成并不矛盾, 二者使用同一组条件概率. 但训练中的真实前缀与生成中的模型前缀可能不同, 错误积累属于分布与使用方式的问题, 不是三角掩码能自动消除的现象.

\begin{exercise}
保持前 5 个输入 token 不变, 替换之后所有 token. 因果模型前 5 个位置的 logits 应怎样变化? 若保持了注意力掩码, 却沿位置轴做 LayerNorm, 这个结论是否仍成立?
\end{exercise}
\begin{solution}
在本节所列因果操作下, 前 5 个位置的 logits 应完全相同, 除浮点实现差异外不应依赖未来. 位置轴归一化把未来输入纳入前缀的均值和方差, 从而可以改变前缀, 结论失效. 这说明因果不变量是在检验信息路径, 不是仅检验注意力模块的形状.
\end{solution}

\originalsection{位置表征与旋转位置编码}
若没有位置表示也没有掩码, 内容自注意力具有排列等变性. 设 $\Pi$ 是把行重新排序的置换矩阵, 则 $Q'=\Pi Q,K'=\Pi K,V'=\Pi V$, 分数为 $S'=\Pi S\Pi^\top$. 这先重排行再在每行重排候选, 逐行 softmax 因而给 $A'=\Pi A\Pi^\top$, 最后
$C'=A'V'=\Pi A\Pi^\top\Pi V=\Pi C$.
交换输入顺序只会交换相应输出, 不能额外表达“原先谁在谁之前”. 因果掩码本身区分可读范围, 不满足任意置换下的这一假设; 不能把未掩码的结论直接用于带固定三角掩码的网络.

一种位置方案是 $X_i=E(x_i)+P_i$, 让位置向量进入所有后续投影. 可以学习有限位置表, 也可以用确定的正弦函数生成向量. 这使模型能够利用位置, 但有限位置表没有给出未训练位置的参数. 正弦可计算任意整数位置, 也不代表训练学到的运算对任意长度都可靠.

旋转位置编码 RoPE 直接改变查询与键的匹配几何\cite{rope}. 为避免行列约定混淆, 先定义作用于二维列向量的旋转矩阵
\[
R(\alpha)=\begin{pmatrix}\cos\alpha&-\sin\alpha\\\sin\alpha&\cos\alpha\end{pmatrix}.
\]
对行向量 $q_i,k_j\in\R^{1\times2}$, 令
$\widetilde q_i=q_iR(i\omega)^\top$,
$\widetilde k_j=k_jR(j\omega)^\top$.
两个二维向量按各自的位置旋转, 然后取点积. 由三角加法公式, $R(\alpha)^\top R(\beta)=R(\beta-\alpha)$, 所以
\[
\widetilde q_i\widetilde k_j^\top
=q_iR(i\omega)^\top R(j\omega)k_j^\top
=q_iR((j-i)\omega)k_j^\top.
\]
结果依赖相对距离 $j-i$, 而不是两个绝对位置各自进入一个任意函数. 这条等式对固定内容向量严格成立. 网络深层的 $q_i,k_j$ 本身已经依赖上下文, 因而整模型的所有行为不只由相对距离决定.

对于偶数维 $d_k$, 把坐标分成 $d_k/2$ 对, 第 $r$ 对采用自己的频率 $\omega_r$, 独立做上述旋转, 总点积是各对点积之和. 在这一简化方案中旋转只作用于查询与键, 值仍保留原投影. 旋转正交, 因而保持每对向量的长度, 但改变不同位置之间的点积. 若某一对 $q_i=k_j=(1,0)$, 它的点积就是 $\cos((j-i)\omega)$, 随距离振荡, 并不单调衰减. 多频率和训练表征可以形成其他匹配形状; 不能把单个旋转恒等式当成“所有注意力随距离单调减弱”的证明.

已知位置和频率固定时, RoPE 的反向仍是线性映射. 从 $\widetilde q_i=q_iR_i^\top$ 得 $G_{q_i}=G_{\widetilde q_i}R_i$, 键同理. 即先对旋转后的查询与键做注意力反向, 再乘逆方向的旋转返回投影输出. 有些扩展训练频率或缩放参数, 那时还需对这些参数求导; 本节没有假定它们可训练.

\begin{exercise}
只用一个二维旋转对, $\omega=\pi/2$, $q_i=k_j=(1,0)$. 计算相对距离 $j-i=0,1,2,3,4$ 的点积. 这能否证明无限上下文外推?
\end{exercise}
\begin{solution}
点积依次为 $1,0,-1,0,1$. 它周期重复, 可以区分部分距离, 也可能对不同距离给相同匹配. RoPE 的公式能在任意位置计算, 但这些数值并不能证明模型对未训练长度的内容、检索和生成仍可靠. 外推能力须通过预先规定的长上下文任务与长度对照评价, 而不是由可计算性推出.
\end{solution}

\originalsection{多头、残差与逐 token 网络}
一个头的输出是一个值混合, 同一组权重同时作用于所有值坐标. 多头让不同投影子空间分别选择混合. 对 $h=1,\ldots,H$, 每头得到 $C^{(h)}\in\R^{T\times d_v}$, 拼接为 $C_*\in\R^{T\times Hd_v}$, 再用 $W_O\in\R^{Hd_v\times d}$ 得 $O=C_*W_O$. 典型设计令 $Hd_k=Hd_v=d$, 但这不是定义所必需. 拼接仅排列特征, 不把不同头的权重平均成一个头.

一个有明确计算路径的 pre-norm 块为
\begin{align*}
U&=\operatorname{LN}_1(X),& H_1&=X+\operatorname{MHA}(U),\\
R&=\operatorname{LN}_2(H_1),&
F&=\phi(RW_F+b_F),&Y&=H_1+FW_G+b_G.
\end{align*}
$W_F\in\R^{d\times r},W_G\in\R^{r\times d}$, $F$ 的形状为 $T\times r$. 两个 LayerNorm 都逐 token 作用, 有各自的参数, 导数见第 \ref{sec:layernorm} 节. 原始 Transformer 的归一化位置与此式不同\cite{attention}; 学习一个实现时应把实际前向写出来, 而不是认为所有叫 Transformer 的模块有同一个图. 本书的小模型也会故意省略归一化, 以便把故障定位到较少对象.

输出分类器令 logits 为 $YW_U+b_U$, $W_U\in\R^{d\times V_0}$, 概率为其 softmax. 多层残差流把初始词表向量、各层注意力信息和局部非线性变换持续组合, 形成预测所用表征. 这条函数图解释了信息能怎样传递, 尚未告诉我们训练实际学出了哪一种算法. “某头在检索前例”或“某方向编码一个属性”是需要设计干预和对照来检验的机制假说, 不能单凭结构名称认定.

一些现代模型使用 RMSNorm: 对单 token, $r=(D^{-1}\sum_i x_i^2+\epsilon)^{-1/2}$, 输出 $y_i=\gamma_i r x_i$. 它不减均值, 因而不具有 LayerNorm 的整体平移不变性. 固定正 $\epsilon$ 时, 令 $u=g\odot\gamma$, 用与上一章相同的微分计算可得
$G_x=ru-r^3x\,\overline{ux}$, $G_\gamma=g\odot(rx)$, 共享参数仍对 token 求和. 若阅读的模型使用它, 不能把 LayerNorm 的两个均值项直接搬过去. 选择哪种归一化改善哪种训练行为, 属于给定架构和训练设置的实验问题.

\originalsection{注意力的完整反向传播}
从下游传来的 $G_C\in\R^{T\times d_v}$ 开始. $C=AV$ 给 $dC=dA\,V+A\,dV$, 由 Frobenius 内积收集系数得
\[
G_A=G_CV^\top\in\R^{T\times T},\qquad
G_V=A^\top G_C\in\R^{T\times d_v}.
\]
这两条路径分别控制混合权重与被混合的信息. 若只更新值而不更新查询和键, 就删掉了通过匹配权重的学习路径.

固定一行, 对允许分数 $s_l$ 求导. 商法则给
$\partial A_j/\partial s_l=A_j(\mathbf1\{j=l\}-A_l)$. 因而 $ds_l$ 的系数为
\[
(G_S)_{il}=A_{il}\left((G_A)_{il}-\sum_j A_{ij}(G_A)_{ij}\right).
\]
它先减去本行权重下的平均贡献, 再乘该位置的权重. 若某位置的损失系数高于该行的加权平均, 提高其分数会使损失的一阶近似增大, 所以对这些分数直接做梯度下降时会倾向降低它; 低于平均的位置相反. 这是由损失系数决定的局部变化, 不是保证训练会发现某个语义位置. 禁止位置权重为零, 分数梯度也为零; 掩码是常量且不更新. 全行梯度和为零, 与分数同时加常数不改变 softmax 相符.

在未禁止的分数上, $dS=(dQ\,K^\top+Q\,dK^\top)/\sqrt{d_k}$, 所以
\[
G_Q=\frac{G_SK}{\sqrt{d_k}}\in\R^{T\times d_k},\qquad
G_K=\frac{G_S^\top Q}{\sqrt{d_k}}\in\R^{T\times d_k}.
\]
然后返回投影参数与共享输入:
\begin{align*}
G_{W_Q}&=X^\top G_Q,&G_{W_K}&=X^\top G_K,&G_{W_V}&=X^\top G_V,\\
G_X&=G_QW_Q^\top+G_KW_K^\top+G_VW_V^\top.
\end{align*}
三条输入路径都属于同一 $X$, 所以相加. 有 RoPE 时先反旋转再代入投影式; 有独立查询输入的交叉注意力时, 查询梯度与键值输入梯度分别保留, 不能都加到同一个矩阵.

多头的输出投影也按同一规则求导: $G_{W_O}=C_*^\top G_O$, $G_{C_*}=G_OW_O^\top$. 按原拼接边界把 $G_{C_*}$ 切为各头的 $T\times d_v$ 梯度, 各头独立使用上述链. 每头有自己的投影参数, 但共享输入, 所以各头的输入梯度再相加. 拼接的反向是切片, 不是求平均.

对 pre-norm 块, 设 $G_Y$ 已知, 从第二条残差开始. 记
$G_F=G_YW_G^\top$,
$D=(G_F\odot\phi'(RW_F+b_F))$,
$G_R=DW_F^\top$.
同时有 $G_{W_G}=F^\top G_Y$, $G_{W_F}=R^\top D$, 两个偏置沿位置求和. 然后
\[
G_{H_1}=G_Y+\operatorname{LN}_2^{\rm back}(G_R),\qquad
G_X=G_{H_1}+\operatorname{LN}_1^{\rm back}
\bigl(\operatorname{MHA}^{\rm back}(G_{H_1})\bigr).
\]
这里 $\operatorname{LN}^{\rm back}$ 表示用该层前向保存的均值、方差和参数执行上一章的反向, 不是再对梯度做一次前向归一化; 其 $\gamma,\beta$ 梯度也须收集. 两个显式加项是恒等残差路径. 这就把 logits 的误差完整送到本块输入, 可继续穿过前面各块, 最后汇总到词表和位置参数.

若输出分类器与词表共享权重, 例如 $Z=YE^\top$, 则 $E$ 还有来自分类器的贡献 $G_E^{\rm out}=G_Z^\top Y$. 查表路径贡献与该项相加才是总梯度. 共享减少参数数目, 不减少它在计算图中出现的路径数. 批量情形中这些矩阵参数梯度均对所有文档和位置求和, 但激活梯度保留其索引.

\originalsection{两位置的一次完整更新}\label{sec:attention-numeric}
公式链容易在转置处失去直觉. 用 $B=1,T=2,D=1$ 把所有量压成可手算的矩阵, 输入 $X=(1,2)^\top$, 三个投影权重均为标量 1. 这里只含因果注意力, 没有前馈层、残差或输出分类器. 标量损失是末位置输出与目标 $1.5$ 的平方误差 $\ell=(C_2-1.5)^2/2$, 第一位置不直接计入损失.

前向有 $Q=K=V=(1,2)^\top$, 未掩码分数为 $\left(\begin{smallmatrix}1&2\\2&4\end{smallmatrix}\right)$. 掩码删除第一行的第二个候选. 记 $a=1/(1+e^2)\approx0.119203$, 则
\[
A=\begin{pmatrix}1&0\\a&1-a\end{pmatrix},\qquad
C=\begin{pmatrix}1\\2-a\end{pmatrix},\qquad
r=C_2-1.5\approx0.380797.
\]
由 $d\ell=r\,dC_2$ 得 $G_C=(0,r)^\top$. 再回到乘法与 softmax:
\[
G_A=\begin{pmatrix}0&0\\r&2r\end{pmatrix},\quad
G_V=\begin{pmatrix}ar\\(1-a)r\end{pmatrix},\quad
G_S=\begin{pmatrix}0&0\\-a(1-a)r&a(1-a)r\end{pmatrix}.
\]
第一行虽影响第一位置输出, 该输出没有下游损失, 所以它的梯度全为零. 第二行的两个分数梯度和为零, 符合 softmax 对整体平移不敏感. 令 $u=a(1-a)r\approx0.0399813$, 有 $G_Q=G_SK=(0,u)^\top$, $G_K=G_S^\top Q=(-2u,2u)^\top$. 所以
\[
G_{W_Q}=X^\top G_Q=2u,\qquad G_{W_K}=X^\top G_K=2u,
\qquad G_{W_V}=X^\top G_V=(2-a)r.
\]
三个数分别约为 $0.0799625,0.0799625,0.716202$. 即使输入不更新, 它对上游模型的梯度仍可算出: $G_X=G_Q+G_K+G_V\approx(-0.0345704,0.455349)^\top$. 若 $X$ 本身由查表产生, 这个量才能继续传回词表与位置参数.

取学习率 $0.1$, 同时用当前梯度更新三个投影, 而固定输入 $X$. 所有更新必须使用同一次前向得到的旧权重与梯度, 不能更新一个权重后再用新值计算另一权重的梯度. 更新与重新前向的实际数值如下; \texttt{attention\_walkthrough} 还输出每个中间矩阵, 并对全部三个参数作中心差分.
更新后的 $(W_Q,W_K,W_V)$ 为 $(0.992004,0.992004,0.928380)$. 重新前向得到 $C_2=1.742951$, 损失从 $0.07250321$ 降至 $0.02951257$. 中心差分的最大误差为 $7.8\times10^{-12}$.


\originalsection{公式与批量实现}\label{sec:attention-code}
第 \ref{sec:tensor} 节已解释查表、广播与展平. 在原单头、无归一化的 \texttt{TinyDecoder} 中, \texttt{q,k,v} 都是 $(B,T,D)$, 每个批次独立做注意力. 分数的定义是 $s_{bij}=\sum_dq_{bid}k_{bjd}/\sqrt D$: 一边保留查询位置 $i$, 另一边保留键位置 $j$, 只收缩特征轴. 因而要把 \texttt{k} 的最后两轴交换成 $(B,D,T)$.

\texttt{k.transpose(0,2,1)} 表示新数组的第 $(b,d,t)$ 项是原数组的第 $(b,t,d)$ 项; 批次轴仍在最前. 三维数组的 \texttt{k.T} 会反转全部轴, 得 $(D,T,B)$, 它不是所需的逐批转置. NumPy 的 \texttt{@} 对最后两个轴执行矩阵乘法, 对前面的批次轴匹配或广播. 如下代码逐行对应数学对象:
\begin{verbatim}
x = E[ids] + P[:T]             # (B,T,D)
q, k, v = (x @ W for W in (WQ, WK, WV))
s = q @ k.transpose(0,2,1) / np.sqrt(D)
s = np.where(mask, s, -np.inf) # mask: (T,T)
a = softmax(s)                 # (B,T,T), last axis
c = a @ v                      # (B,T,D)
h = x + c @ WO                 # attention residual
f = np.tanh(h @ WF + bF)        # (B,T,2D)
z = h + f @ WG + bG            # FFN residual
logits = z @ WU + bU           # (B,T,V0)
\end{verbatim}
这里 \texttt{mask} 的形状 $(T,T)$ 沿 $B$ 广播, 每行至少保留自身. \texttt{softmax} 沿最后一轴归一化, 所以固定查询位置与批次后, 在允许的键位置之间分配权重. 表中的代码变量 \texttt{WQ} 等在实际程序中存于 \texttt{p['Q']} 等参数字典, \texttt{dl} 是平均交叉熵对 \texttt{logits} 的梯度.

反向的对象、形状与关键代码可以从输出逐级核对:
\begin{center}
\begin{tabular}{lll}\toprule
对象 & 形状 & \texttt{loss\_grad} 中的运算 \\\midrule
$G_U$ & $D\times V_0$ & \texttt{mm(z,dl)} \\
$G_z$ & $B\times T\times D$ & \texttt{dl@p['U'].T} \\
$G_F$ & $D\times 2D$ & \texttt{mm(h,df)} \\
$G_O$ & $D\times D$ & \texttt{mm(c,dh)} \\
$G_A$ & $B\times T\times T$ & \texttt{dc@v.transpose(0,2,1)} \\
$G_V$ & $B\times T\times D$ & \texttt{a.transpose(0,2,1)@dc} \\
$G_Q$ & $B\times T\times D$ & \texttt{ds@k/np.sqrt(self.d)} \\
$G_{W_Q}$ & $D\times D$ & \texttt{mm(x,dq)} \\
$G_P$ & $T\times D$ & \texttt{dx.sum(0)} \\\bottomrule
\end{tabular}
\end{center}
矩阵参数的 $G_U$ 指输出投影权重, 与上节单头 $G_V$ 的值激活梯度不要混淆. \texttt{df} 已包含 $\tanh$ 的导数, \texttt{dh} 已把 $z=h+fW_G+b_G$ 的直接路径和前馈路径相加. 对 $h=x+CW_O$ 又有一条直接路径, 所以
\[
G_x=G_h+G_QW_Q^\top+G_KW_K^\top+G_VW_V^\top.
\]
这一行的 $G_h$ 不能漏掉. 矩阵梯度按批次与位置求和, 激活梯度则保留它们; 只记形状而不记共享关系仍会写错.

\begin{exercise}
运行 \texttt{residual\_debug} 实验. 它只在反向中漏掉上述 $G_h$ 项, 前向完全不变. 先预测哪些参数组的有限差分会失败, 再根据输出定位, 最后给出修复. 为什么这比只看损失是否下降更有诊断力?
\end{exercise}
\begin{solution}
故障在回到输入 $x$ 时发生, 因此继续依赖 $G_x$ 的词表 $E$ 与位置表 $P$ 会失败; 此前已算出的注意力投影、前馈层与输出层参数梯度仍正确. 修复是把 \texttt{dh} 加回 \texttt{dx}. 程序用 $B=2,T=3,D=4$, 词表 6, 种子 19, 对 206 个参数坐标全部检查. 本次结果如下:
正确实现的最大误差为 $3.8\times10^{-11}$; 漏项后 $E,P$ 的误差分别为 $0.120850$ 与 $0.177792$, 其余参数组仍小于 $10^{-7}$.

漏梯度仍可能产生下降方向, 所以损失变小不能确认反向实现. 有限差分按参数组定位可将问题缩到共享输入的上游, 再通过计算图确定漏掉的恒等路径. 它仍是此输入与初始化上的数值诊断, 不是对所有情形的正确性证明.
\end{solution}

\originalsection{计算代价与缓存}
注意力把每个查询与每个允许键匹配. 一般稠密情形, $QK^\top$ 需要 $O(T^2d_k)$ 运算, $AV$ 需要 $O(T^2d_v)$. $H$ 个头的匹配和聚合共为 $O(HT^2(d_k+d_v))$, 若 $Hd_k=Hd_v=d$, 就是 $O(T^2d)$. 投影和输出映射还有约 $O(Td^2)$ 的开销, 前馈层另有 $O(Tdr)$. 因果三角只把位置对数量从 $T^2$ 改为约 $T(T+1)/2$, 没有把数量级变成线性.

若显式保存每头的分数或概率, 存储需 $O(HT^2)$ 个数; 查询、键、值和输出本身约需 $O(Td)$. 训练还要保存各层用于反向的激活. 不能把只算一层推理的内存当成完整训练内存. 运算次数、峰值存储和内存搬运是不同指标, 一个实现减少存储, 不意味着减少了全部乘加次数.

逐步生成时, 旧位置的键和值可缓存. 新查询只需与此前的键匹配, 然后读相应值, 单步这部分代价约为 $O(T H(d_k+d_v))$, 加上投影和前馈代价. 每层 KV 缓存需 $O(T H(d_k+d_v))$ 个数. 这避免对完整前缀反复计算旧键值, 但随着已生成长度增长, 新 token 仍要读取越来越长的缓存. 还必须保持位置编号、掩码和训练时定义一致. 本章的 CPU 实验没有实现缓存性能基准, 上述是由形状和使用次数得出的计算分析.

\originalsection{分块在线 softmax 的公式与证明}
显式存储 $T\times T$ 分数不是计算精确注意力的数学要求. 困难在于一个查询的 softmax 分母依赖所有允许键: 不能给每个块单独归一化后直接拼接, 否则块间概率质量错误. 我们需要在处理键块时保存可合并的统计量, 让最终输出和一次处理所有键相同.

先固定一个查询. 对已处理的非空允许键集合 $\mathcal J$, 记
\[
m=\max_{j\in\mathcal J}s_j,\qquad
l=\sum_{j\in\mathcal J}e^{s_j-m},\qquad
u=\sum_{j\in\mathcal J}e^{s_j-m}v_j.
\]
$m,l$ 是标量, $u$ 是 $d_v$ 维行向量. 当前输出为 $u/l$. 保存未归一化的 $u$ 是为了下一块到来时还能改变共同分母. 对与已处理集合不相交的新非空允许块 $\mathcal B$（$\mathcal J\cap\mathcal B=\varnothing$）, 取新最大值
$m'=\max(m,\max_{j\in\mathcal B}s_j)$, 更新
\begin{align*}
l'&=e^{m-m'}l+\sum_{j\in\mathcal B}e^{s_j-m'},\\
u'&=e^{m-m'}u+\sum_{j\in\mathcal B}e^{s_j-m'}v_j.
\end{align*}
每次只使用不大于零的指数参数, 避免原始大分数直接指数上溢. 旧最大值变化时, 旧累积量须一起缩放, 不能只缩放新块.

\begin{proposition}
在上述不相交条件下, 更新后 $m',l',u'$ 恰为并集 $\mathcal J\cup\mathcal B$ 的相应统计量. 在有限实数精确算术中, 处理完所有允许键后 $u/l$ 等于稠密 softmax 注意力输出.
\end{proposition}
\begin{proof}
最大值的式子显然是并集最大值. 对旧键 $j\in\mathcal J$, 有
$e^{s_j-m'}=e^{m-m'}e^{s_j-m}$. 因而旧键在新共同尺度下的指数和是 $e^{m-m'}l$, 加上新块就得到并集指数和 $l'$. 对每项再乘向量 $v_j$, 相同恒等式给 $u'$. 以首个非空允许块为初始统计量, 逐块归纳即可. 最后分子与分母同乘 $e^{-m}$ 不改变比值, 故 $u/l=\sum_j e^{s_j}v_j/\sum_j e^{s_j}$.
\end{proof}
也可以使用 $m=-\infty,l=0,u=0$ 的空状态, 但实现必须特别处理还没有允许键的行, 避免计算 $-\infty-(-\infty)$. 全禁止块应跳过, 首个非空块再建立有限最大值. 每个查询最终至少有一个允许键, 才能除以 $l>0$. 这在因果和 padding 组合掩码中尤其重要.

\begin{example}
分数为 $(0,\log3)$, 值为 $(2,0),(0,4)$, 每个键单独成为一个块. 用在线公式复算输出.
\end{example}
\begin{solution}
处理第一个键后 $m=0,l=1,u=(2,0)$. 第二个键使 $m'=\log3$, 旧尺度因子为 $1/3$, 新键指数为 1. 所以 $l'=1/3+1=4/3$, $u'=(2/3,4)$, 输出 $u'/l'=(1/2,3)$, 与稠密结果相同. 如果两个块分别归一化, 它们都给权重 1, 便会把总输出错算为 $(2,4)$.
\end{solution}

实际分块还把查询分为至多 $b_q$ 行的小块, 每次与至多 $b_k$ 个键相乘. 一个分数块形状仅为 $b_q\times b_k$, 各行分别维护 $m,l,u$. 除输入输出外, 当前块及其累积量约需 $O(b_qb_k+b_qd_v)$ 工作空间. 固定块大小时, 不必存全部 $T^2$ 个分数; 但所有稠密查询--键对仍需计算, 算术数量级仍为二次. 若实现另外存着一个完整 $T\times T$ 布尔掩码, 这份掩码本身仍占二次空间; 想让总额外空间也避免二次增长, 还须按索引在块内产生因果掩码, 或用合适的稀疏表示.

反向也可以重算概率块. 前向保存每行 $z_i=m_i+\log l_i$, 则允许位置 $A_{ij}=e^{s_{ij}-z_i}$. 给定 $G_C$, 记
\[
D_i=G_{C_i}C_i^\top.
\]
它恰等于 softmax 反向所需的行平均贡献, 因为
\[
\sum_j A_{ij}(G_A)_{ij}
=\sum_jA_{ij}G_{C_i}V_j^\top
=G_{C_i}\left(\sum_j A_{ij}V_j\right)^\top
=G_{C_i}C_i^\top.
\]
因此在一个重算的块上, 计算 $(G_A)_{ij}=G_{C_i}V_j^\top$, 再算 $(G_S)_{ij}=A_{ij}((G_A)_{ij}-D_i)$; 按已有公式累积 $G_Q,G_K,G_V$. $D_i$ 无须等待整行 $G_A$ 显式存好才求出. 这是避免保存完整概率矩阵的一条数学路径, 代价是反向重算一些前向量.

FlashAttention 将分块、在线归一化与重算结合, 专门管理 GPU 高带宽内存和片上存储之间的数据搬运\cite{flash}. 本节证明了一个核心代数不变量, 没有证明某个 CPU 循环达到论文的 GPU IO 复杂度或速度. 减少分数工作空间、实现融合内核和取得实际加速是不同层次的工作. 原论文的实现、硬件和预算需要另行评价.

\originalsection{完整标准 Transformer 的 CPU 配对训练}\label{sec:standard-model}
手算展示一条注意力路径，完整模型还要把位置、归一化、前馈、残差和读出连接起来。\path{standard_transformer_study.py} 使用 PyTorch 的标准层实现两层前置归一化的因果解码器；它不是包装外部 API，也没有下载预训练权重。结构如下：词表 10、宽度 $32$、两头、前馈宽度 $64$、上下文长度 7，两个残差块，GELU 激活，最后归一化及线性读出，总计 18,026 个可训练参数。位置是可学习嵌入，没有使用本章另一节推导的 RoPE；实现和方法名不能混淆。

每一块可写成
\[
U=X+\operatorname{MHA}(\operatorname{LN}(X)),\qquad
X'=U+\operatorname{FFN}(\operatorname{LN}(U)),
\]
其中 $\operatorname{FFN}(u)=\operatorname{GELU}(uW_1+b_1)W_2+b_2$。训练对整个词表用交叉熵，不只监督叫声。八条原创语法句从说话者、颜色、旁观者各两种组合构成；固定六条训练、两条留出，整个设计在训练前写成独立 JSON 并计算哈希。

我们比较均衡六条训练句与将同一条训练句额外重复 30 次。若单条损失是 $\ell_i$，第二组目标为
\[
L_{\rm repeat}=\frac{31\ell_0+\sum_{i\ne0}\ell_i}{36},
\qquad L_{\rm balanced}=\frac16\sum_{i=0}^5\ell_i.
\]
唯一句子数没有增加，分布权重改变了。两组同一种子从相同参数出发、做 600 次 AdamW 更新，但每步分别处理 6 与 36 行，因此没有匹配 FLOPs。这个实验检验频率重加权及其伴随处理成本，不能当成等算力数据质量消融，更不是 FineWeb 的大规模复现。

种子固定为 144、145、146，学习率 $0.003$、权重衰减 $0.01$、梯度裁剪 1，无 dropout，CPU 单线程 float32。六次运行全部报告，未用留出成绩挑最好种子。下面由实际结果 JSON 自动生成，代码与环境哈希另附。
\begin{center}\begin{tabular}{clrrr}\toprule
种子 & 条件 & 加权训练 NLL & 均衡训练 NLL & 留出 NLL \\ \midrule
144 & 均衡 & 0.2565 & 0.2565 & 1.3159 \\
144 & 单条重复 & 0.0899 & 0.4310 & 1.3747 \\
145 & 均衡 & 0.2565 & 0.2565 & 1.2958 \\
145 & 单条重复 & 0.0900 & 0.4311 & 1.3475 \\
146 & 均衡 & 0.2565 & 0.2565 & 1.2773 \\
146 & 单条重复 & 0.0900 & 0.4311 & 1.3985 \\
\bottomrule\end{tabular}\end{center}

重复条件的自身加权训练损失更低，三个配对留出损失却更高。若只看优化器监控的训练目标，会错把偏向常见句的改进当作总体泛化提高；使用同一均衡训练分布评价又得到不同结果。这里两条留出样本太少，不能从三次初始化推出广泛的数据结论。

六个模型对原始、改变旁观者及改变说话者三种叫声测试均符合该有限语法；但这些干预组合与训练句子重叠，不能识别唯一内部机制，也不能证明任意组合迁移。编辑未来 token 后先前位置的最大 logits 差为零；额外测试还检查了未来嵌入的梯度为零。这提供了因果掩码实现证据。想继续研究，应事先扩充组合、独立留出角色和词表，匹配 token 处理预算，再比较数据重加权；不能看到当前结果后悄悄修改这一次已冻结实验。

\originalsection{实验、机制假说与练习}
新增 \texttt{attention} 实验在 CPU 的 NumPy float64 中, 用 $T=37,d_k=5,d_v=7$ 的随机输入和因果掩码, 比较稠密计算与查询--键块大小 $1\times1,4\times6,8\times7,50\times50$ 的在线计算. 固定种子 144 时, 最大绝对输出差不超过 $4.45\times10^{-16}$. 把查询与键各放大 100 倍, 使分数进入强饱和区, 两种算法也保持有限且一致. 这检验数值等价, 不把浮点计算声称为实数精确相等.

同一实验用共享输入产生查询、键和值, 替换第 15 位置及之后的输入, 因果掩码下前 15 个位置的输出最大变化为 0; 取消掩码后的最大变化约为 37.1055. 注意力反向还对一个任意线性下游损失的全部 68 个 $Q,K,V$ 坐标进行中心差分, 最大混合误差低于 $8.87\times10^{-11}$. 其他检验包括禁止位置权重为零、每行和为 1、首块全禁止时仍正确处理后续允许块、全禁止查询被明确拒绝. 这些输入特意覆盖常见实现边界, 仍不是对所有程序状态的形式证明.

默认 $8\times7$ 块的单个分数数组最多含 56 个数, 而稠密分数数组有 $37^2=1369$ 个数. 这两个值来自实际配置和数组形状, 不是进程峰值内存测量. 程序为教学方便仍传入完整布尔掩码, 输出和中间累积量也占内存. 环境记录没有 GPU 速度或端到端训练吞吐量; 不应把这里的工作数组比较写成 FlashAttention 性能复现.

这些结构如何连接基础模型能力? 注意力提供按内容与位置读取上下文的路径, 前馈层变换当前表征, 残差允许多个层次的信息叠加. 预训练让这些路径共同服务后继预测. 在上下文学习任务中, 提示中的例子可以作为当次输入被读取, 无须在推理时更新参数. 但“有路径”不等于“训练一定学会利用它”. 要研究某个模型是否从示例学到对应规则, 可固定模型, 对照有无示例、示例标签随机化、示例顺序变化和未见组合, 再检查目标答案的条件概率与实际正确率. 它们区分对样例的读取、模板捷径与任务规律, 不能由一个注意力热力图直接确定.

\begin{exercise}
$X$ 形状为 $12\times32$, 每头 $W_Q,W_K$ 为 $32\times8$, $W_V$ 为 $32\times16$, 一共有 4 个头. 写出每头和拼接后的形状, 以及输出回到 32 维所需 $W_O$ 的形状.
\end{exercise}
\begin{solution}
每头 $Q,K$ 为 $12\times8$, $V,C$ 为 $12\times16$, $S,A$ 为 $12\times12$. 拼接后 $C_*$ 为 $12\times64$, 所以 $W_O$ 为 $64\times32$, 输出为 $12\times32$. 给定输出梯度, $G_{C_*}=G_OW_O^\top$ 也是 $12\times64$, 按四个 16 列块切分后分别进行各头反传. 各头的输入梯度都是 $12\times32$, 最后相加.
\end{solution}
\begin{exercise}
一个查询有三个键, 分数为 $(\log2,0,\log3)$, 标量值为 $(1,4,2)$. 第一块含前两个键, 第二块含最后一个键. 求在线状态与最终输出.
\end{exercise}
\begin{solution}
第一块最大值 $m=\log2$, 指数为 $(1,1/2)$, 故 $l=3/2,u=1+4/2=3$. 新最大值为 $m'=\log3$, 旧缩放因子为 $2/3$. 所以 $l'=(2/3)(3/2)+1=2$, $u'=(2/3)3+2=4$, 输出为 2. 稠密式也给 $(2\cdot1+1\cdot4+3\cdot2)/(2+1+3)=2$.
\end{solution}
\begin{exercise}
若某个查询仅允许一个键, 它的查询和键分数梯度是否为零? 值梯度呢? 最大注意力权重能否直接表示最大因果作用?
\end{exercise}
\begin{solution}
该行权重恒为 1, softmax 对唯一分数的导数为 $1(1-1)=0$, 因而该行对查询和键的分数路径不产生梯度. 但输出等于那个值, 值仍接收 $G_C$. 键也可能被其他查询使用, 所以整个键向量的总梯度不一定为零. 最大权重还要乘值并经过输出投影和其他残差路径, 因而不能单独确定因果作用. 可设计保持其他因素可比的输入替换或中间激活替换, 检查指定输出变化, 并核对替换是否带来分布外伪影.
\end{solution}
\begin{exercise}
若前向已存 $C_i$ 和 $z_i$, 为什么反向能够逐块求 $G_S$ 而不先保存整行 $G_A$? 写出必须仍累积的参数或激活梯度.
\end{exercise}
\begin{solution}
因为整行的加权平均 $\sum_jA_{ij}(G_A)_{ij}$ 已等于可直接计算的 $D_i=G_{C_i}C_i^\top$. 每块重算 $A_{ij}=e^{s_{ij}-z_i}$, 用 $A_{ij}((G_A)_{ij}-D_i)$ 得该块 $G_S$. 每块贡献必须累积到查询梯度 $G_Q$, 键梯度 $G_K$ 和值梯度 $G_V$, 再汇总到共享投影矩阵. 查询跨键块重复, 键和值跨查询块重复, 任一方向覆盖而不相加都会漏路径. 减少存储没有消除共享梯度的求和要求.
\end{solution}
\label{ch:end:07}
